% Options for packages loaded elsewhere
\PassOptionsToPackage{unicode}{hyperref}
\PassOptionsToPackage{hyphens}{url}
\documentclass[
  ignorenonframetext,
  brazilian,
  aspectratio=169,
  10pt]{beamer}
\newif\ifbibliography
\usepackage{pgfpages}
\setbeamertemplate{caption}[numbered]
\setbeamertemplate{caption label separator}{: }
\setbeamercolor{caption name}{fg=normal text.fg}
\beamertemplatenavigationsymbolsempty
% remove section numbering
\setbeamertemplate{part page}{
  \centering
  \begin{beamercolorbox}[sep=16pt,center]{part title}
    \usebeamerfont{part title}\insertpart\par
  \end{beamercolorbox}
}
\setbeamertemplate{section page}{
  \centering
  \begin{beamercolorbox}[sep=12pt,center]{section title}
    \usebeamerfont{section title}\insertsection\par
  \end{beamercolorbox}
}
\setbeamertemplate{subsection page}{
  \centering
  \begin{beamercolorbox}[sep=8pt,center]{subsection title}
    \usebeamerfont{subsection title}\insertsubsection\par
  \end{beamercolorbox}
}
% Prevent slide breaks in the middle of a paragraph
\widowpenalties 1 10000
\raggedbottom
\AtBeginPart{
  \frame{\partpage}
}
\AtBeginSection{
  \ifbibliography
  \else
    \frame{\sectionpage}
  \fi
}
\AtBeginSubsection{
  \frame{\subsectionpage}
}
\usepackage{iftex}
\ifPDFTeX
  \usepackage[T1]{fontenc}
  \usepackage[utf8]{inputenc}
  \usepackage{textcomp} % provide euro and other symbols
\else % if luatex or xetex
  \usepackage{unicode-math} % this also loads fontspec
  \defaultfontfeatures{Scale=MatchLowercase}
  \defaultfontfeatures[\rmfamily]{Ligatures=TeX,Scale=1}
\fi
\usepackage{lmodern}
\usetheme[]{Madrid}
\usefonttheme[]{professionalfonts}
\ifPDFTeX\else
  % xetex/luatex font selection
\fi
% Use upquote if available, for straight quotes in verbatim environments
\IfFileExists{upquote.sty}{\usepackage{upquote}}{}
\IfFileExists{microtype.sty}{% use microtype if available
  \usepackage[]{microtype}
  \UseMicrotypeSet[protrusion]{basicmath} % disable protrusion for tt fonts
}{}
\makeatletter
\@ifundefined{KOMAClassName}{% if non-KOMA class
  \IfFileExists{parskip.sty}{%
    \usepackage{parskip}
  }{% else
    \setlength{\parindent}{0pt}
    \setlength{\parskip}{6pt plus 2pt minus 1pt}}
}{% if KOMA class
  \KOMAoptions{parskip=half}}
\makeatother
\ifLuaTeX
\usepackage[bidi=basic,shorthands=off,]{babel}
\else
\usepackage[bidi=default,shorthands=off,]{babel}
\fi
\ifLuaTeX
  \usepackage{selnolig} % disable illegal ligatures
\fi
\setlength{\emergencystretch}{3em} % prevent overfull lines
\providecommand{\tightlist}{%
  \setlength{\itemsep}{0pt}\setlength{\parskip}{0pt}}
\usepackage{amsmath}
\usepackage{amssymb}
\usepackage{booktabs}
\usepackage{xcolor}
\definecolor{AzulAula}{HTML}{1D4ED8}
\definecolor{AzulEscuro}{HTML}{0F172A}
\definecolor{CinzaTexto}{HTML}{334155}
\definecolor{VerdeAula}{HTML}{047857}
\setbeamercolor{structure}{fg=AzulAula}
\setbeamercolor{title}{fg=white,bg=AzulEscuro}
\setbeamercolor{frametitle}{fg=white,bg=AzulEscuro}
\setbeamercolor{normal text}{fg=CinzaTexto,bg=white}
\setbeamercolor{block title}{fg=white,bg=AzulAula}
\setbeamercolor{block body}{fg=AzulEscuro,bg=blue!5}
\setbeamercolor{alertblock title}{fg=white,bg=red!75!black}
\setbeamercolor{alertblock body}{fg=AzulEscuro,bg=red!5}
\setbeamertemplate{navigation symbols}{}
\setbeamertemplate{caption}[numbered]
\setbeamertemplate{footline}[frame number]
\setbeamerfont{caption}{size=\scriptsize}
\setlength{\abovecaptionskip}{2pt}
\setlength{\belowcaptionskip}{0pt}
\renewcommand{\tablename}{Tabela}
\renewcommand{\figurename}{Figura}
\usepackage{bookmark}
\IfFileExists{xurl.sty}{\usepackage{xurl}}{} % add URL line breaks if available
\urlstyle{same}
\hypersetup{
  pdftitle={Aula 7: propriedades do estimador de MQO},
  pdfauthor={FLS 6183 - Métodos Quantitativos de Pesquisa II},
  pdflang={pt-BR},
  hidelinks,
  pdfcreator={LaTeX via pandoc}}

\title{Aula 7: propriedades do estimador de MQO}
\subtitle{Distribuição amostral, viés, consistência, Gauss--Markov e
bootstrap}
\author{FLS 6183 - Métodos Quantitativos de Pesquisa II}
\date{14 de outubro de 2026}

\begin{document}
\frame{\titlepage}

\begin{frame}{BLP e MQO}
\protect\phantomsection\label{blp-e-mqo}
Na população, \(\mathbf x=(1,X_1,\ldots,X_k)'\) e
\(\boldsymbol\beta=(\alpha,\beta_1,\ldots,\beta_k)'\): \[
Y=\mathbf x'\boldsymbol\beta+\varepsilon,
\qquad E[\mathbf x\varepsilon]=\mathbf 0,
\qquad
\boldsymbol\beta=E[\mathbf x\mathbf x']^{-1}E[\mathbf xY].
\]

Na amostra, \(\mathbf X\) tem as linhas \(\mathbf x_i'\) e
\(\mathbf Y=(Y_1,\ldots,Y_n)'\): \[
\widehat{\boldsymbol\beta}=(\mathbf X'\mathbf X)^{-1}\mathbf X'\mathbf Y,
\qquad
\widehat{\mathbf e}=\mathbf Y-\mathbf X\widehat{\boldsymbol\beta}.
\]

\begin{block}{Estimador e estimativa}
$\widehat{\boldsymbol\beta}$ é uma função da amostra aleatória.
Cada amostra observada produz uma estimativa dos mesmos coeficientes populacionais.
\end{block}
\end{frame}

\begin{frame}{Nova amostra, nova reta}
\protect\phantomsection\label{nova-amostra-nova-reta}
População conhecida: \(X\sim N(10,3^2)\), \(\varepsilon\sim N(0,2^2)\)
independente de \(X\), \(Y=5+1{,}5X+\varepsilon\).

\begin{figure}

{\centering \includegraphics[width=0.88\linewidth]{../output/aula_07/figuras/01_duas_amostras} 

}

\caption{Duas amostras independentes de 40 pares (X,Y) da mesma população. A linha contínua é o MQO da amostra e a tracejada é o BLP populacional. Fonte: dados sintéticos do processo descrito.}\label{fig:duas-amostras}
\end{figure}

\small

Inclinações estimadas: \(1{,}514\) e \(1{,}552\). O coeficiente
populacional é \(\beta=1{,}5\).
\end{frame}

\begin{frame}{Distribuição amostral}
\protect\phantomsection\label{distribuiuxe7uxe3o-amostral}
Repetimos a amostragem e o ajuste de MQO 5.000 vezes, sempre com
\(n=40\).

\begin{figure}

{\centering \includegraphics[width=0.88\linewidth]{../output/aula_07/figuras/02_distribuicao_amostral} 

}

\caption{Histograma em escala de densidade de 5.000 inclinações de MQO. Cada valor vem de uma amostra independente de tamanho 40. A linha tracejada marca beta = 1,5. Fonte: simulação da população conhecida.}\label{fig:distribuicao-amostral}
\end{figure}

\small

A média das estimativas, \(1{,}502\), aproxima \(E[\widehat\beta]\). O
desvio-padrão, \(0{,}107\), aproxima a dispersão amostral.
\end{frame}

\begin{frame}{Amostras pequenas e não viés}
\protect\phantomsection\label{amostras-pequenas-e-nuxe3o-viuxe9s}
Mesmo processo gerador, com \(\beta=1{,}5\). Repetimos o MQO 5.000 vezes
para cada \(n\).

\begin{figure}

{\centering \includegraphics[width=0.84\linewidth]{../output/aula_07/figuras/07_amostras_pequenas_n10_n40}

}

\caption{Distribuições de 5.000 inclinações por tamanho de amostra, com eixos e largura de classe comuns. A linha tracejada marca o parâmetro 1,5 e a pontilhada, a média das estimativas. Fonte: simulação da mesma população conhecida.}\label{fig:amostras-pequenas}
\end{figure}

\small

Médias: \(1{,}503\) (\(n=10\)) e \(1{,}502\) (\(n=40\)). Desvios-padrão:
\(0{,}254\) e \(0{,}107\).

Com \(n=10\), há maior dispersão individual e a média das repetições
fica próxima de \(1{,}5\).
\end{frame}

\begin{frame}{Viés do estimador}
\protect\phantomsection\label{viuxe9s-do-estimador}
Para explicitar o tamanho da amostra, escrevemos \(\widehat\beta_n\): \[
\boxed{\operatorname{Viés}_n(\widehat\beta_n)=E[\widehat\beta_n]-\beta.}
\] A esperança considera todas as amostras possíveis de tamanho \(n\) da
mesma população.

\begin{block}{Não viés para um tamanho de amostra fixo}
$E[\widehat\beta_n]=\beta$. A definição usa a distribuição amostral do estimador
e não exige tomar o limite quando $n$ cresce.
\end{block}

\begin{block}{Erro de uma estimativa observada}
Na primeira amostra, $\widehat\beta-\beta=
1{,}514-1{,}5$.
Essa diferença é o erro de estimação daquela amostra.
O viés descreve sua média na distribuição amostral.
\end{block}
\end{frame}

\begin{frame}{Variância e erro quadrático médio}
\protect\phantomsection\label{variuxe2ncia-e-erro-quadruxe1tico-muxe9dio}
A variância amostral descreve a dispersão do estimador: \[
\operatorname{V}(\widehat\beta_n)
=E\!\left[(\widehat\beta_n-E[\widehat\beta_n])^2\right].
\]

O erro quadrático médio (EQM) considera a distância ao parâmetro: \[
\boxed{\operatorname{EQM}(\widehat\beta_n)
=E[(\widehat\beta_n-\beta)^2]
=\operatorname{V}(\widehat\beta_n)+\operatorname{Viés}_n(\widehat\beta_n)^2.}
\]

\begin{block}{Não viés e precisão}
Quando o viés é zero, o EQM é igual à variância amostral.
Um estimador não viesado ainda pode produzir estimativas muito dispersas.
\end{block}

\small

O desvio-padrão da distribuição amostral,
\(\sqrt{\operatorname{V}(\widehat\beta_n)}\), é o erro-padrão do
estimador.
\end{frame}

\begin{frame}{Consistência}
\protect\phantomsection\label{consistuxeancia}
\(\widehat\beta_n\) é consistente para \(\beta\) se, para todo
\(\delta>0\), \[
\boxed{\Pr\bigl(|\widehat\beta_n-\beta|>\delta\bigr)\longrightarrow 0
\quad\text{quando } n\longrightarrow\infty.}
\] Escrevemos \(\widehat\beta_n\xrightarrow{p}\beta\): convergência em
probabilidade.

\begin{block}{O papel de $n$}
Consistência descreve uma sequência de estimadores conforme a amostra cresce.
Não viés descreve a esperança do estimador para um $n$ fixo.
\end{block}

\small

Exemplo com \(Y_i\) i.i.d., \(E[Y_i]=\mu\neq0\) e
\(0<\operatorname{V}(Y_i)<\infty\):

\begin{center}
\begin{tabular}{lll}
\toprule
Estimador de $\mu$ & Não viesado & Consistente \\
\midrule
$Y_1$, qualquer que seja $n$ & Sim & Não \\
$\dfrac{n}{n+1}\bar Y$ & Não & Sim \\
\bottomrule
\end{tabular}
\end{center}
\end{frame}

\begin{frame}{Calcular MQO na amostra}
\protect\phantomsection\label{calcular-mqo-na-amostra}
Com dados observados finitos, o MQO minimiza \[
S(\mathbf b)=(\mathbf Y-\mathbf X\mathbf b)'(\mathbf Y-\mathbf X\mathbf b).
\]

\begin{block}{Condição para coeficientes únicos}
$\mathbf X$ tem posto completo: suas $k+1$ colunas são linearmente independentes.
Assim, $\mathbf X'\mathbf X$ tem inversa e
\[
\widehat{\boldsymbol\beta}=(\mathbf X'\mathbf X)^{-1}\mathbf X'\mathbf Y.
\]
\end{block}

\begin{block}{Independência não é exigida para calcular}
Podemos calcular o MQO com observações dependentes ou de distribuições distintas.
Segundos momentos populacionais finitos também não são uma exigência desse cálculo.
\end{block}

\small

Sem posto completo, o problema de minimização tem soluções, mas os
coeficientes não são únicos.
\end{frame}

\begin{frame}{MQO: reescrever a fórmula conhecida}
\protect\phantomsection\label{mqo-reescrever-a-fuxf3rmula-conhecida}
Começamos com a fórmula que já estudamos: \[
\widehat{\boldsymbol\beta}=(\mathbf X'\mathbf X)^{-1}\mathbf X'\mathbf Y.
\]

A lei dos grandes números trata da convergência de médias. Vamos
escrever o MQO em termos dessas médias.

Dividir uma matriz por \(n\) multiplica sua inversa por \(n\): \[
\left(\frac{\mathbf X'\mathbf X}{n}\right)^{-1}
=n(\mathbf X'\mathbf X)^{-1}.
\] Por isso, dividir os dois produtos por \(n\) preserva o resultado: \[
\begin{aligned}
\left(\frac{\mathbf X'\mathbf X}{n}\right)^{-1}
\left(\frac{\mathbf X'\mathbf Y}{n}\right)
&=n(\mathbf X'\mathbf X)^{-1}\frac{\mathbf X'\mathbf Y}{n}\\
&=(\mathbf X'\mathbf X)^{-1}\mathbf X'\mathbf Y
=\widehat{\boldsymbol\beta}.
\end{aligned}
\]

\small

A reescrita conserva exatamente a estimativa calculada na mesma amostra.
\end{frame}

\begin{frame}{Os produtos matriciais contêm somas}
\protect\phantomsection\label{os-produtos-matriciais-contuxeam-somas}
Com um único preditor, cada linha de \(\mathbf X\) é \((1,X_i)\). A
multiplicação de matrizes produz \[
\mathbf X'\mathbf X=
\begin{pmatrix}
n & \sum_i X_i\\
\sum_i X_i & \sum_i X_i^2
\end{pmatrix},
\qquad
\mathbf X'\mathbf Y=
\begin{pmatrix}\sum_i Y_i\\\sum_i X_iY_i\end{pmatrix}.
\] Dividindo por \(n\), cada soma se torna uma média. \(\overline{X^2}\)
é a média dos quadrados e \(\overline{XY}\) a média dos produtos: \[
\frac{\mathbf X'\mathbf X}{n}=
\begin{pmatrix}1 & \bar X\\\bar X & \overline{X^2}\end{pmatrix},
\qquad
\frac{\mathbf X'\mathbf Y}{n}=
\begin{pmatrix}\bar Y\\\overline{XY}\end{pmatrix}.
\]

\begin{block}{Ligação com a fórmula da inclinação}
São as mesmas médias que aparecem na fórmula escalar:
\[
\widehat\beta_1=
\frac{\overline{XY}-\bar X\bar Y}{\overline{X^2}-\bar X^2}
=\frac{\widehat{\operatorname{Cov}}(X,Y)}{\widehat{\operatorname{V}}(X)}.
\]
\end{block}
\end{frame}

\begin{frame}{Consistência: médias e alvo populacional}
\protect\phantomsection\label{consistuxeancia-muxe9dias-e-alvo-populacional}
Cada entrada dos produtos divididos por \(n\) é uma média amostral. Com
vários preditores, a mesma lógica vale.

\begin{block}{Suposições da demonstração}
\begin{itemize}
\item Amostragem i.i.d.: os pares $(Y_i,\mathbf x_i)$ são independentes e identicamente distribuídos.
\item Segundos momentos finitos de $Y$ e de cada preditor.
\end{itemize}
\end{block}

Pela lei dos grandes números, \[
\frac{\mathbf X'\mathbf X}{n}\xrightarrow{p}E[\mathbf x\mathbf x'],
\qquad
\frac{\mathbf X'\mathbf Y}{n}\xrightarrow{p}E[\mathbf xY].
\] Se \(E[\mathbf x\mathbf x']\) tem inversa, a continuidade da inversão
e da multiplicação implica \[
\widehat{\boldsymbol\beta}
=\left(\frac{\mathbf X'\mathbf X}{n}\right)^{-1}
\left(\frac{\mathbf X'\mathbf Y}{n}\right)
\xrightarrow{p}
\underbrace{E[\mathbf x\mathbf x']^{-1}E[\mathbf xY]}_{\boldsymbol\beta:\ \text{coeficientes do BLP}}.
\]
\end{frame}

\begin{frame}{Média condicional do erro}
\protect\phantomsection\label{muxe9dia-condicional-do-erro}
\begin{columns}[T,onlytextwidth]
\begin{column}{0.47\textwidth}
\begin{block}{Propriedade do erro do BLP}
\[
E[\mathbf x\varepsilon]=\mathbf 0.
\]
\end{block}
\end{column}
\begin{column}{0.49\textwidth}
\begin{block}{Suposição chave para o não viés}
\[
E[\varepsilon\mid\mathbf x]=0.
\]
\end{block}
\end{column}
\end{columns}

\begin{block}{Linear nos parâmetros}
Na equação $Y=\mathbf x'\boldsymbol\beta+\varepsilon$, a suposição implica
$E[Y\mid\mathbf x]=\mathbf x'\boldsymbol\beta$.

Por exemplo, com $\mathbf x=(1,X,X^2)'$,
\[
E[Y\mid X]=\alpha+\beta_1X+\beta_2X^2.
\]
A linearidade é em $\alpha,\beta_1,\beta_2$.
A relação com $X$ pode ter curvatura.
\end{block}

\small

Média condicional zero implica ortogonalidade, pela lei da esperança
iterada. A implicação inversa não vale em geral.
\end{frame}

\begin{frame}{Exemplo político: gasto de campanha e votos}
\protect\phantomsection\label{exemplo-poluxedtico-gasto-de-campanha-e-votos}
Considere candidatos a prefeito. Queremos interpretar \(\tau\) como
efeito do gasto: \[
Y_i=\alpha+\tau G_i+\varepsilon_i.
\] \(Y_i\) é a porcentagem de votos; \(G_i\) é o gasto de campanha em
R\$ por eleitor; \(\varepsilon_i\) reúne fatores não medidos, como a
força eleitoral prévia.

\begin{block}{O que exige independência na média?}
Entre candidatos que gastam R\$ 2, R\$ 5 ou R\$ 10 por eleitor, a média de $\varepsilon$
precisaria ser a mesma:
\[
E[\varepsilon\mid G=2]=E[\varepsilon\mid G=5]
=E[\varepsilon\mid G=10]=E[\varepsilon].
\]
\end{block}

Candidatos previamente mais fortes podem atrair mais doações e gastar
mais. Isso pode fazer a média de \(\varepsilon\) variar com \(G\).

\small

Centrar o erro para que \(E[\varepsilon]=0\) é uma normalização. Só com
independência na média obtemos \(E[\varepsilon\mid G]=0\).

\scriptsize

Exemplo hipotético adaptado da lógica de Cunningham,
\textit{Causal Inference: The Mixtape}, seção 2.12. Aqui \(\varepsilon\)
representa fatores não medidos na equação causal proposta; a condição de
média zero não decorre da ortogonalidade do BLP.
\end{frame}

\begin{frame}{Não viés do MQO}
\protect\phantomsection\label{nuxe3o-viuxe9s-do-mqo}
Reunimos os erros do modelo em
\(\boldsymbol\varepsilon=(\varepsilon_1,\ldots,\varepsilon_n)'\).

\begin{block}{Suposição chave para o não viés}
Sob amostragem i.i.d., $E[\varepsilon_i\mid\mathbf x_i]=0$ implica
$\boxed{E[\boldsymbol\varepsilon\mid\mathbf X]=\mathbf 0.}$
\end{block}

Na fórmula de MQO, substituímos \(\mathbf Y\) por
\(\mathbf X\boldsymbol\beta+\boldsymbol\varepsilon\) e aplicamos a
distributiva: \[
\begin{aligned}
\widehat{\boldsymbol\beta}
&=(\mathbf X'\mathbf X)^{-1}\mathbf X'\mathbf Y
=(\mathbf X'\mathbf X)^{-1}\mathbf X'(\mathbf X\boldsymbol\beta+\boldsymbol\varepsilon)\\
&=(\mathbf X'\mathbf X)^{-1}(\mathbf X'\mathbf X)\boldsymbol\beta
+(\mathbf X'\mathbf X)^{-1}\mathbf X'\boldsymbol\varepsilon\\
&=\boldsymbol\beta+(\mathbf X'\mathbf X)^{-1}\mathbf X'\boldsymbol\varepsilon.
\end{aligned}
\] Ao condicionar em \(\mathbf X\), a matriz dos preditores fica fixa:
\[
E[\widehat{\boldsymbol\beta}\mid\mathbf X]
=\boldsymbol\beta+(\mathbf X'\mathbf X)^{-1}\mathbf X'
\underbrace{E[\boldsymbol\varepsilon\mid\mathbf X]}_{\mathbf 0}
=\boldsymbol\beta.
\]

\small

Com posto completo e esperança do estimador existente, a lei da
esperança iterada dá \(E[\widehat{\boldsymbol\beta}]=\boldsymbol\beta\).
O resultado vale para \(n\) fixo.
\end{frame}

\begin{frame}{Variância do MQO}
\protect\phantomsection\label{variuxe2ncia-do-mqo}
Suponha
\(\operatorname{V}(\boldsymbol\varepsilon\mid\mathbf X)=\sigma^2\mathbf I_n\).
Então: \[
\boxed{\operatorname{V}(\widehat{\boldsymbol\beta}\mid\mathbf X)
=\sigma^2(\mathbf X'\mathbf X)^{-1}.}
\]

\small

Com um preditor e intercepto,
\(\operatorname{V}(\widehat\beta\mid\mathbf X)=
\sigma^2/\sum_i(X_i-\bar X)^2\). Maior ruído aumenta a variância. Maior
variação do preditor reduz a variância.
\end{frame}

\begin{frame}{Teorema de Gauss--Markov}
\protect\phantomsection\label{teorema-de-gaussmarkov}
Considere
\(\mathbf Y=\mathbf X\boldsymbol\beta+\boldsymbol\varepsilon\), com
posto completo, \[
E[\boldsymbol\varepsilon\mid\mathbf X]=\mathbf 0,
\qquad
\operatorname{V}(\boldsymbol\varepsilon\mid\mathbf X)=\sigma^2\mathbf I_n.
\]

\begin{block}{Versão clássica}
O MQO tem a menor variância entre os estimadores lineares em $\mathbf Y$
e não viesados condicionalmente em $\mathbf X$.
\end{block}

\begin{itemize}
\item Linear em $\mathbf Y$: $\widehat{\boldsymbol\beta}^{\,a}=\mathbf A\mathbf Y$,
com $\mathbf A$ dependendo apenas de $\mathbf X$.
\item Não viés para todo $\boldsymbol\beta$:
$E[\widehat{\boldsymbol\beta}^{\,a}\mid\mathbf X]=\boldsymbol\beta$.
\item A comparação vale para cada coeficiente e para qualquer combinação linear dos coeficientes.
\end{itemize}

\small

BLUE: \textit{best linear unbiased estimator}, ou melhor estimador
linear não viesado. O teorema dispensa a normalidade dos erros.
\end{frame}

\begin{frame}{Outra prova do não viés do MQO}
\protect\phantomsection\label{outra-prova-do-nuxe3o-viuxe9s-do-mqo}
Defina \(\mathbf A_0=(\mathbf X'\mathbf X)^{-1}\mathbf X'\). Como
\(\mathbf A_0\) depende apenas de \(\mathbf X\),
\(\widehat{\boldsymbol\beta}=\mathbf A_0\mathbf Y\). Com
\(\mathbf Y=\mathbf X\boldsymbol\beta+\boldsymbol\varepsilon\) e
\(E[\boldsymbol\varepsilon\mid\mathbf X]=\mathbf 0\): \[
\begin{aligned}
E[\widehat{\boldsymbol\beta}\mid\mathbf X]
&=E[\mathbf A_0\mathbf Y\mid\mathbf X]\\
&=\mathbf A_0 E[\mathbf Y\mid\mathbf X]\\
&=\mathbf A_0(\mathbf X\boldsymbol\beta
+E[\boldsymbol\varepsilon\mid\mathbf X])\\
&=\mathbf A_0\mathbf X\boldsymbol\beta.
\end{aligned}
\]

\begin{block}{Condição de não viés}
Para que essa média seja $\boldsymbol\beta$ para todo vetor de coeficientes,
precisamos de $\mathbf A_0\mathbf X=\mathbf I_{k+1}$:
a identidade deixa $\boldsymbol\beta$ inalterado.
\end{block}

Como \(\mathbf X'\mathbf X\) é invertível: \[
\mathbf A_0\mathbf X=(\mathbf X'\mathbf X)^{-1}(\mathbf X'\mathbf X)
=\mathbf I_{k+1}.
\]

\small

\(\mathbf I_{k+1}\) tem uma linha e uma coluna por coeficiente,
incluindo o intercepto. A conclusão usa média condicional zero; não
exige variância constante.
\end{frame}

\begin{frame}{Comparação entre estimadores}
\protect\phantomsection\label{comparauxe7uxe3o-entre-estimadores}
Agora considere outro estimador linear
\(\widehat{\boldsymbol\beta}^{\,a}
=\mathbf A\mathbf Y\), com \(\mathbf A\) dependendo apenas de
\(\mathbf X\). A mesma conta dá
\(E[\widehat{\boldsymbol\beta}^{\,a}\mid\mathbf X]
=\mathbf A\mathbf X\boldsymbol\beta\). Para não viés para todo
\(\boldsymbol\beta\), \(\mathbf A\mathbf X=\mathbf I_{k+1}\).

Escreva \(\mathbf A=\mathbf A_0+\mathbf C\). Como
\(\mathbf A\mathbf X=\mathbf A_0\mathbf X=\mathbf I_{k+1}\), segue
\(\mathbf C\mathbf X=\mathbf 0\). Transpondo e usando
\(\mathbf A_0=(\mathbf X'\mathbf X)^{-1}\mathbf X'\): \[
\mathbf X'\mathbf C'=(\mathbf C\mathbf X)'=\mathbf 0,
\qquad
\mathbf A_0\mathbf C'=(\mathbf X'\mathbf X)^{-1}\mathbf X'\mathbf C'=\mathbf 0.
\] \[
\begin{aligned}
\operatorname{V}(\mathbf A\mathbf Y\mid\mathbf X)
&=\sigma^2(\mathbf A_0+\mathbf C)(\mathbf A_0+\mathbf C)'\\
&=\sigma^2\mathbf A_0\mathbf A_0'+\sigma^2\mathbf C\mathbf C'.
\end{aligned}
\]

\begin{block}{A parcela adicional tem variância não negativa}
Para qualquer vetor $\mathbf a$ de pesos dos coeficientes,
\[
\operatorname{V}(\mathbf a'\mathbf A\mathbf Y\mid\mathbf X)
-\operatorname{V}(\mathbf a'\widehat{\boldsymbol\beta}\mid\mathbf X)
=\sigma^2(\mathbf C'\mathbf a)'(\mathbf C'\mathbf a)\geq0.
\]
\end{block}
\end{frame}

\begin{frame}{Monte Carlo: tamanho da amostra}
\protect\phantomsection\label{monte-carlo-tamanho-da-amostra}
Mesmo processo gerador: \(Y=5+1{,}5X+\varepsilon\), com \(\sigma=2\).

\begin{figure}

{\centering \includegraphics[width=0.88\linewidth]{../output/aula_07/figuras/03_tamanho_amostra} 

}

\caption{Distribuições de inclinações em 5.000 amostras independentes por tamanho n. Mesmos eixos e mesma largura de classe nos três painéis. A linha tracejada marca beta = 1,5. Fonte: simulação com X normal e erro independente.}\label{fig:tamanho-amostra}
\end{figure}

\small

Desvios-padrão: \(0{,}107\) (\(n=40\)), \(0{,}053\) (\(n=160\)) e
\(0{,}026\) (\(n=640\)). A redução da dispersão ilustra a consistência
demonstrada anteriormente.
\end{frame}

\begin{frame}{Parâmetro da regra e coeficiente do BLP}
\protect\phantomsection\label{paruxe2metro-da-regra-e-coeficiente-do-blp}
Agora, \(Y=5+1{,}5X+u\), com \(u=0{,}8(X-10)+v\) e \(E[v\mid X]=0\). \[
E[u\mid X]=0{,}8(X-10),\qquad
Y=-3+2{,}3X+v.
\]

\begin{figure}

{\centering \includegraphics[width=0.84\linewidth]{../output/aula_07/figuras/05_parametro_e_blp} 

}

\caption{Inclinações em 5.000 amostras por n da nova regra geradora. A linha pontilhada marca 1,5 e a tracejada marca o coeficiente do BLP, 2,3. Fonte: dados sintéticos, com X normal e v independente.}\label{fig:parametro-blp}
\end{figure}

\small

O MQO é consistente para o BLP, cujo coeficiente é \(2{,}3\). \(u\) é o
erro da regra original. O erro da projeção linear é \(\varepsilon=v\).
Recuperar o parâmetro \(1{,}5\) da regra original exige condições
adicionais.
\end{frame}

\begin{frame}{Monte Carlo e bootstrap}
\protect\phantomsection\label{monte-carlo-e-bootstrap}
\begin{columns}[T,onlytextwidth]
\begin{column}{0.47\textwidth}
\begin{block}{Monte Carlo}
Conhecemos a população e a regra geradora.

\medskip
Geramos novas amostras independentes e calculamos o estimador em cada uma.

\medskip
Podemos comparar as estimativas com o parâmetro conhecido.
\end{block}
\end{column}
\begin{column}{0.49\textwidth}
\begin{block}{Bootstrap}
Temos uma amostra observada da população.

\medskip
Reamostramos suas unidades e calculamos o estimador em cada reamostra.

\medskip
Usamos a distribuição empírica como aproximação da população.
\end{block}
\end{column}
\end{columns}

\bigskip

O bootstrap aplica a lógica do plug-in à distribuição amostral do
estimador.
\end{frame}

\begin{frame}{Bootstrap de pares}
\protect\phantomsection\label{bootstrap-de-pares}
Cada unidade é um par \((X_i,Y_i)\). Para uma amostra de \(n\) unidades:

\begin{enumerate}
\item Sorteamos $n$ pares da amostra, com reposição.
\item Ajustamos o MQO e guardamos a inclinação $\widehat\beta^{*(b)}$.
\item Repetimos o procedimento $B$ vezes.
\item Calculamos o desvio-padrão das $B$ inclinações.
\end{enumerate}

\begin{block}{Estimativa bootstrap do erro-padrão}
\[
\widehat{\operatorname{EP}}_{\mathrm{boot}}(\widehat\beta)
=\sqrt{\frac{1}{B-1}\sum_{b=1}^{B}
\left(\widehat\beta^{*(b)}-\overline{\widehat\beta^{*}}\right)^2},
\qquad
\overline{\widehat\beta^{*}}=\frac1B\sum_{b=1}^{B}\widehat\beta^{*(b)}.
\]
\end{block}

\small

O asterisco identifica uma estimativa em reamostra. O bootstrap de pares
usual pressupõe unidades i.i.d. e condições de regularidade. O sorteio
mantém \(X_i\) e \(Y_i\) da mesma unidade juntos.
\end{frame}

\begin{frame}{Distribuição bootstrap}
\protect\phantomsection\label{distribuiuxe7uxe3o-bootstrap}
Usamos a primeira amostra de 40 pares, com \(\widehat\beta=
1{,}514\), e fazemos \(B=5.000\) reamostras.

\begin{figure}

{\centering \includegraphics[width=0.88\linewidth]{../output/aula_07/figuras/06_bootstrap} 

}

\caption{Monte Carlo com 5.000 novas amostras (esquerda) e bootstrap com 5.000 reamostras de uma única amostra (direita). A linha tracejada marca beta = 1,5 à esquerda e a estimativa original à direita. Fonte: dados sintéticos.}\label{fig:distribuicao-bootstrap}
\end{figure}

\small

Dispersão Monte Carlo: \(0{,}107\). Erro-padrão bootstrap: \(0{,}084\).
O bootstrap aproxima a incerteza usando a amostra disponível. O
erro-padrão bootstrap também é uma estimativa e varia entre amostras.
\end{frame}

\begin{frame}{Síntese}
\protect\phantomsection\label{suxedntese}
\begin{itemize}
\item O viés compara a esperança do estimador com o parâmetro.
O EQM combina viés ao quadrado e variância amostral.
\item A consistência descreve a aproximação ao parâmetro conforme $n$ cresce.
Para o BLP, amostragem i.i.d., segundos momentos finitos e matriz populacional invertível
formam um conjunto de hipóteses suficientes para a consistência.
\item A média condicional zero permite demonstrar o não viés do MQO.
Sob as hipóteses da versão clássica de Gauss–Markov, o MQO tem a menor variância
entre estimadores lineares não viesados.
\item O bootstrap usa a distribuição empírica para aproximar a distribuição amostral
e estimar o erro-padrão.
\end{itemize}

\bigskip

Na próxima aula: análise de resíduos e diagnóstico do modelo.
\end{frame}

\begin{frame}{Leituras}
\protect\phantomsection\label{leituras}
\begin{block}{Leituras previstas no programa}
Aronow, Peter M., e Benjamin T. Miller. 2019.
\textit{Foundations of Agnostic Statistics}. Cambridge University Press.
Capítulo 3, seção 3.2.2: teoria da estimação, pp. 102--105.

\medskip
Hansen, Bruce E. 2022. \textit{Econometrics}. Princeton University Press.
Capítulo 4, seções 4.4--4.8: modelo linear, esperança e variância do MQO,
momentos incondicionais e Gauss--Markov.
Capítulo 7, seção 7.2: consistência do MQO.
\end{block}

\begin{block}{Complemento para a introdução ao bootstrap}
Aronow e Miller, capítulo 3, seção 3.4.3, pp. 130--132.
\end{block}

\scriptsize

As figuras usam dados sintéticos. As aproximações de Monte Carlo e
bootstrap usam 5.000 repetições por cenário. Os resultados teóricos
dependem das condições indicadas nos slides.
\end{frame}

\end{document}
